\documentclass{64110}

\author{Emira Ibrahimović}
\problem{4}
\collab{Claude helped me format the \LaTeX file}

\begin{document}
\begin{enumerate}

% ---------------------------------------------------------------
\item \textbf{Fire Problem}

\begin{enumerate}

\item[(i)]
The components of a state are:
\begin{itemize}
  \item $M \times N$ possible robot locations
  \item $2$ possible values for \texttt{carried\_patient}
  \item $2^{M \times N}$ possible values for the state of the fire grid
\end{itemize}
This gives a total state space of size $M \times N \cdot 2^{MN+1}$, which is enormous.
If we approached this with dynamic programming, at every relevant $t$ we would need to
compute the optimal policy for each of these states, which is infeasible.

\item[(ii)]
We compute $P(F_{ij}^{(t+1)} \mid F^{(t)})$.
Let $F_{\mathrm{neighs}}$ be the random sub-matrix of $F^{(t)}$ containing the
neighbors of cell $(i,j)$ and the cell itself.
Let $F_r$ range over $\{0,1\}^{|F_{\mathrm{neighs}}|}$, the set of all such
realized matrices.
Let $K_{i'j'} = W[i'-i+1,\; j'-j+1]$.

Using the formulas given in the problem:
\begin{align*}
P(F_{ij}^{(t+1)} \mid F^{(t)})
&= P(F_{ij}^{(t+1)} \mid F_{\mathrm{neighs}}) \\
&= \sum_{F_r} P(F_{\mathrm{neighs}} = F_r)
   \sum_{i',j'} K_{i'j'} (F_r)_{i'j'} \\
&= \sum_{i',j'} K_{i'j'}
   \sum_{F_r} P(F_{\mathrm{neighs}} = F_r)\,(F_r)_{i'j'}.
\end{align*}
Since $(F_r)_{i'j'} \in \{0,1\}$, the inner sum equals $P(F^{(t)}_{\text{neighbor }(i',j')\text{ of }(i,j)} = 1) = F^{(t)}_{i'j'}$.
Therefore:
\[
  P(F_{ij}^{(t+1)} \mid F_{\mathrm{neighs}}) = \sum_{i',j'} K_{i'j'} F^{(t)}_{i'j'}
  = (W * F^{(t)})_{ij},
\]
where $*$ denotes element-wise convolution. Hence
\[
  P(F^{(t+1)} \mid F^{(t)}) = W \circledast P(F^{(t)}).
\]

\item[(iii)]
Given an initial state $s_0$, let $\pi = (s_1, s_2, \ldots, s_T)$ denote a path.
We want:
\begin{align*}
\pi^* &= \arg\max_\pi \prod_{i=1}^{T} P(\text{robot survives step } s_{i-1} \to s_i) \\
      &= \arg\max_\pi \sum_{i=1}^{T} \log P(\text{step } i \text{ is safe}) \\
      &= \arg\max_\pi \sum_{i=1}^{T} \log\bigl(1 - P(\text{onfire}_{s_i})\bigr) \\
      &= \arg\min_\pi \sum_{i=1}^{T} -\log\bigl(1 - P(\text{onfire}_{s_i})\bigr) \\
      &= \arg\min_\pi \sum_{i=1}^{T} c - \log\bigl(1 - P(\text{onfire}_{s_i})\bigr).
\end{align*}
We add a constant $c$ so that the cost is never zero: when $P(\text{onfire}_{s_i}) = 0$,
$\log(1) = 0$. Step costs of 0 make our search algorithms break.

\item[(iv)]
When the agent sees that its path to the hospital is obstructed by fire, it starts
pacing between two squares. Only after the path is unblocked does it continue moving
towards the hospital.

At each time step the agent searches the space of possible futures to find one where it
reaches the goal. It does this by simulating the fire probabilistically -- evolving the
fire probability grid after every time step. When the probabilities of fire along the
chosen path are high, the path cost explodes since $-\log(1-p) \to \infty$ as $p \to 1$.
The agent then prefers paths through safe squares and waits, circling until it exhausts
its step budget.

\end{enumerate}


% ---------------------------------------------------------------
\item \textbf{Stochastic Fire Problem}

\begin{enumerate}

\item[(i)]
The agent chooses to go right and manoeuvre around the fire. Introducing stochasticity
removes the certainty about where fire exists---especially further into the planning
horizon. We also lose certainty about cells that definitely contain no fire. This makes
the agent more cautious and more likely to choose longer but safer paths over shorter
direct ones.

\item[(ii)]
A closed-loop agent would not make a long-term plan based solely on hypothetical rollouts
of the system. Although our A* agent is not entirely open-loop -- before taking a step, it
makes an observation and replans -- it still produces a long-term plan by simulating
the system (in the sense of propagating probabilities through time) without using the actual future
state. The agent's belief can lead it to consider some paths as completely unviable
because they have a high \emph{probability} of being dangerous, whereas the actual
realisation of the fire is such that those paths are viable and even shorter.

\item[(iii)]
If planning were very expensive, I would prefer not to replan after every step. Similarly
to the previous mini-project, we could follow an initial plan and observe the environment as we
go. At each observation we compute the probability that the next planned step leads the
agent into fire. If this probability exceeds a threshold (e.g.\ $5\%$), we replan.
The trade-off is that by following the original plan blindly we could end up surrounded
by fire; when we finally detect this and replan, there may be no remaining viable path.

\item[(iv)]
Most-likely-outcome determinization would be problematic. Consider the following initial
state (H = hospital, \texttt{\#} = fire, \texttt{a} = agent carrying patient,
\texttt{-} = wall):

\begin{verbatim}
. H #
# . #
    -
# . a
# . .
\end{verbatim}

The agent must choose between going left and going down (a wall blocks the upward route).
The cell to the left has four fire neighbors, so the most likely outcome is that it
remains clear in the next step---but with only about $44\%$ probability, which is hardly
guaranteed. The shorter path goes left, so the determinized agent chooses it.
However, moving down first and waiting to see how the fire develops may be the safer
strategy before committing to a route.

\end{enumerate}


% ---------------------------------------------------------------
\item \textbf{MCTS}

\begin{enumerate}

\item[(i)]
Both agents choose the wide, spacious route in order to avoid potential fire on their
path. MCTS is more of a closed-loop agent than A*: it does not plan far into the future
but instead tries to predict as many possible futures as possible and acts according to
the best estimated reward for each action. In some sense MCTS distinguishes different
possible futures better than A* does. A*, by propagating fire probabilities, ``mushes
together'' many potential futures and does not maintain a clear picture of what any
single one looks like.

\item[(ii)]
The advantage of A* lies precisely in this mushing together. MCTS only knows about the
futures it encounters during rollouts: with fewer rollouts it knows less about the ways
the MDP can evolve and makes less informed decisions. A*, by contrast, encodes
information about \emph{all} possible futures through the per-cell fire probabilities.

In particular, MCTS might roll out roughly equal numbers of futures starting with action
\textit{right} and action \textit{down}. However, many more rollouts of action
\textit{right} are needed to get a reliable estimate of its value. Given a large rollout
budget, MCTS will start realising that \textit{right} yields higher reward and will
prioritise it. With small budgets the outcomes are too noisy to reveal this.

\item[(iii)]
Depending on how dangerous the surrounding area actually is, the state value might not
be affected very much. With $n_{\mathrm{simulations}} = 10$, a single bad rollout from a
given state incurs a large negative reward, but this can be compensated by the positive
rewards of the remaining nine rollouts. If the area is genuinely dangerous, it is more
likely that all rollouts lead to bad outcomes, and the state value reflects that
correctly.

\end{enumerate}

\end{enumerate}
\end{document}